\section{第10组：热力学第一定律}
\sourceproblem{10-T1}{60}{64}
\problemtext 从a经b至c吸热200 J且对外做功80 J；从a经d至c吸热144 J；从c回a外界对系统做功52 J，分别求内能差、第二路径功及返回热量。
\workings 取吸热 $Q$ 和系统对外做功 $W$ 为正，$\Delta U=Q-W$。$U_c-U_a=120$ J，故c内能较大；第二路径 $W=144-120=24$ J；返回时 $\Delta U=-120$ J、$W=-52$ J，故 $Q=-172$ J，即放热172 J。
\crosscheck 内能为状态函数，功和热量依赖路径；回路净热等于净功。
\sourceproblem{10-T2}{60}{64}
\problemtext 初压 $1.0\times10^5$ Pa、初体积 $1.0\times10^{-3}$ 立方米的氧气，从0摄氏度升至100摄氏度，分别求等压与等容热量、做功。
\workings $\nu R=p_0V_0/T_0=100/273.15$ J/K，$\Delta T=100$ K。刚性双原子理想模型中
\[
Q_V=\frac52\nu R\Delta T\approx91.5\,\mathrm J,\quad W_V=0,\quad
Q_p=\frac72\nu R\Delta T\approx128.1\,\mathrm J,\quad W_p=\nu R\Delta T\approx36.6\,\mathrm J.
\]
\crosscheck 两过程内能增量相同；$Q_p-Q_V=W_p$，采用理想气体、定热容近似。
\sourceproblem{10-T3}{61}{65}
\problemtext A经ABC至C吸热326 J、对外做功126 J；从C沿CA返回，外界做功52 J，求返回过程热量。
\workings $U_C-U_A=326-126=200$ J。返回时 $\Delta U=-200$ J，系统对外功为 $-52$ J，故 $Q=-252$ J，即放热252 J。
\crosscheck 完整循环净热 $326-252=74$ J，与净功 $126-52=74$ J相同。
\sourceproblem{10-T4}{61}{65}
\problemtext 1 mol氢气初态 $p_0=10^5$ Pa、$T_0=293.15$ K、$V_0$，末态 $T_1=353.15$ K、$2V_0$。比较先等容升温后等温膨胀与先等温膨胀后等容升温的功、热量和内能变化，并给出p--V图。
\workings 刚性双原子模型 $\Delta U=(5/2)R(60)\approx1247$ J，两路径相同。高温等温膨胀 $W_1=RT_1\ln2\approx2035$ J，$Q_1\approx3282$ J；低温等温膨胀 $W_2=RT_0\ln2\approx1689$ J，$Q_2\approx2936$ J。图可按无量纲坐标 $x=V/V_0,y=p/p_0$ 作出：两路径分别为 $(1,1)\to(1,T_1/T_0)$ 后沿 $y=(T_1/T_0)/x$ 至 $x=2$；以及沿 $y=1/x$ 至 $(2,1/2)$ 后竖直至 $(2,T_1/(2T_0))$。
\begin{center}
\setlength{\unitlength}{0.8mm}
\begin{picture}(92,67)
\put(8,8){\vector(1,0){78}}\put(8,8){\vector(0,1){55}}
\put(84,2){$V$}\put(1,61){$p$}
\put(29,40){\circle*{2}}\put(18,40){$A$}
\put(29,50){\circle*{2}}\put(18,51){$B$}
\put(69,24){\circle*{2}}\put(72,23){$D$}
\put(69,29){\circle*{2}}\put(72,30){$C$}
\put(29,40){\vector(0,1){10}}\put(69,24){\vector(0,1){5}}
\qbezier(29,50)(45,34)(69,29)
\qbezier(29,40)(45,28)(69,24)
\put(36,54){$T_1$}\put(47,21){$T_0$}
\put(25,2){$V_0$}\put(64,2){$2V_0$}
\end{picture}
\end{center}
\crosscheck 示意图不按精确比例绘制，曲线解析式给出准确图形；高温膨胀功较大，热量差等于功差。
\sourceproblem{10-T5}{62}{66}
\problemtext 图1的bca为绝热压缩，b1a与b2a为图示任意压缩路径，判断功与吸热符号；图2同压两平衡态A、B且 $V_B>V_A$，判断任何连接过程必然满足的性质。
\workings 图1三路径均压缩，$W_1,W_2<0$；上方b2a压缩功绝对值较大，下方b1a较小。绝热路径 $Q_c=0$ 给 $\Delta U=-W_c$，故 $Q_1=W_1-W_c>0$、$Q_2=W_2-W_c<0$，选B。图2由 $T=pV/(\nu R)$ 得 $T_B>T_A$，固定组成理想气体内能增加，选B。
\crosscheck 第二问只约束端点等压，并非过程恒压，不能保证功或热量符号。
\sourceproblem{10-T6}{62}{66}
\problemtext 理想气体循环a--b等温、b--c绝热、c--d等压、d--a等容；表中给 $W_{ab}=50$ J、$\Delta U_{bc}=-50$ J、$W_{cd}=-50$ J、$Q_{cd}=-150$ J，填表求热效率。
\workings 由 $Q=\Delta U+W$：ab为 $(\Delta U,W,Q)=(0,50,50)$ J；bc为 $(-50,50,0)$ J；cd为 $(-100,-50,-150)$ J；da因循环内能变化和为零且等容，得 $(150,0,150)$ J。净功50 J、吸收热量200 J，故 $\eta=50/200=25\%$。
\crosscheck 各段热量代数和50 J等于净功；热效率分母为所有正吸热之和而非净吸热。

\sourceproblem{10-T7}{63}{67}
\problemtext 理想气体作两等压、两绝热组成的顺时针循环，B、C温度400和300 K，求热效率。
\workings 采用准静态可逆绝热和恒定热容模型，$p_A=p_B=p_H,\ p_C=p_D=p_L$，且 $T_C/T_B=T_D/T_A=(p_L/p_H)^{(\gamma-1)/\gamma}=3/4$。吸热 $Q_H=\nu C_p(T_B-T_A)$，放热量 $Q_L=\nu C_p(T_C-T_D)$，故
\[
\eta=1-\frac{T_C-T_D}{T_B-T_A}=1-\frac34=25\%.
\]
\crosscheck 不应把B、C直接视为两个恒温热源；相同温度比来自可逆绝热关系，且作为热机须有 $p_H>p_L>0$、$T_B>T_A>0$，从而 $Q_H>0$。仅给一般不可逆绝热过程或任意变热容气体时，数据不足以保证此值。
\sourceproblem{10-T8}{63}{67}
\problemtext 热机在1000 K与300 K热源间工作，分别将高温源升至1100 K或将低温源降至200 K，求理论效率增量。
\workings 卡诺上限 $\eta=1-T_L/T_H$。原效率70\%；第一种为 $1-300/1100=72.727\%$，提高2.727个百分点；第二种为80\%，提高10个百分点。
\crosscheck 两种调整分别与原状态比较，不连续实施；若问相对增长率，则分别为3.90\%和14.29\%。
\sourceproblem{10-T9}{64}{68}
\problemtext 理想制冷机电动机输出200 W，冷室270 K、外界300 K，求10 min内从冷室取热量。
\workings 可逆制冷系数 $\epsilon=T_L/(T_H-T_L)=9$。输入功 $W=200\times600=1.20\times10^5$ J，故 $Q_L=\epsilon W=1.08\times10^6$ J。
\crosscheck 向热源放热 $Q_H=Q_L+W=1.20\times10^6$ J；制冷系数可大于1，不是热机效率。
\sourceproblem{10-T10}{64}{68}
\problemtext 2 mol理想气体300 K可逆等温膨胀，体积由0.02增至0.04立方米，求熵变。
\workings $\Delta S=\int\delta Q_{\rm rev}/T=\nu R\ln(V_2/V_1)=2R\ln2\approx11.52\,\mathrm{J/K}$。
\crosscheck 膨胀熵增为正；系统内能不变不代表熵不变。
\sourceproblem{10-T11}{64}{68}
\problemtext 4.00 mol理想气体体积加倍，分别求400 K等温膨胀功、该等温过程熵变，以及可逆绝热膨胀熵变。
\workings 若等温过程准静态可逆，$W=\nu RT\ln2\approx9.22\times10^3$ J。两等温平衡态间的熵变 $\Delta S=\nu R\ln2\approx23.04\,\mathrm{J/K}$。可逆绝热过程 $\delta Q_{\rm rev}=0$，故 $\Delta S=0$。
\crosscheck 原题第一问仅明确等温，功依实际路径及外压而定；给出的功是教材可逆模型值，不能用于自由膨胀。第二问熵为状态函数，第三问末温通常低于初温。
\sourceproblem{10-T12}{65}{69}
\problemtext 卡诺热机在300 K与100 K间正循环，每循环净功6000 J，画T--S图并求吸热、放热及效率。
\workings $\eta=1-100/300=2/3$，$Q_H=6000/\eta=9000$ J，$Q_L=Q_H-6000=3000$ J。可逆等温吸热熵增 $\Delta S=9000/300=30\,\mathrm{J/K}$；两绝热段等熵。
\begin{center}
\begin{tikzpicture}[>=Stealth,scale=0.85]
\draw[->] (0,0)--(6,0) node[right] {$S$};
\draw[->] (0,0)--(0,4.4) node[above] {$T$};
\draw[->,thick] (1,3.5)--(4.5,3.5);
\draw[->,thick] (4.5,3.5)--(4.5,1);
\draw[->,thick] (4.5,1)--(1,1);
\draw[->,thick] (1,1)--(1,3.5);
\node[left] at (0,3.5) {$300\,\mathrm K$};
\node[left] at (0,1) {$100\,\mathrm K$};
\node[below] at (1,0) {$S_1$};
\node[below] at (4.5,0) {$S_1+30\,\mathrm{J/K}$};
\end{tikzpicture}
\end{center}
\crosscheck T--S矩形面积 $(300-100)\times30=6000$ J；每循环工作物质熵变为零。
\sourceproblem{10-T13}{65}{69}
\problemtext 0.5 kg、0摄氏度冰置于恒温20摄氏度的大热源，最终变为20摄氏度水；求刚融完和随后升温阶段的熵变。熔化热335 kJ/kg，水比热4180 J/(kg K)。
\workings 取 $T_m=273.15$ K、$T_H=293.15$ K。用可逆参考路径，
\[
\Delta S_1=\frac{m\lambda}{T_m}\approx613.22\,\mathrm{J/K},\quad
\Delta S_2=mc\ln\frac{T_H}{T_m}\approx147.69\,\mathrm{J/K}.
\]
总物质熵增约760.90 J/K。
\crosscheck 真实过程存在有限温差而不可逆；计算熵差仍可用同端点可逆路径。第二阶段不应简单用热量除以20摄氏度。
\sourceproblem{10-T14}{65}{69}
\problemtext 接上题，求热源熵变及冰水与热源整体的总熵变。
\workings 恒温热源放热 $Q=m\lambda+mc(20)=209300$ J，故 $\Delta S_H=-209300/293.15\approx-713.97$ J/K。整体熵变
\[
\Delta S_{\rm tot}=\frac{m\lambda}{T_m}+mc\ln\frac{T_H}{T_m}-\frac Q{T_H}\approx46.93\,\mathrm{J/K}.
\]
\crosscheck 总熵增为正，符合不可逆传热；热源温度保持恒定是大热源近似。
